% ===========================================================================
\section{The channel as a conditional law}\label{sec:channel-model}
\warmth{0}\quad\whisper{A channel is a table of conditional probabilities; capacity is the best use we can make of it.}

\begin{strand}
Everything about a discrete channel is contained in one stochastic matrix. The
matrix is given to us by physics; what remains free is the distribution with
which we feed the inputs. Capacity is the value of that freedom.
\end{strand}

\trigger{Reach for this model whenever the corruption of a symbol depends only on that symbol.}

\block{The data of a channel}

\begin{definition}[discrete channel]\label{def:discrete-channel}
A \keyword{discrete channel} is a system consisting of an input alphabet $\X$,
an output alphabet $\Y$, and a probability transition matrix
$p(y\mid x)$, which gives the probability of observing the output symbol $y$
when the input symbol $x$ is sent. Being a stochastic matrix means
\[
  p(y\mid x)\ge0
  \quad\text{for all }x,y,
  \qquad
  \sum_{y\in\Y}p(y\mid x)=\answerin[1.2cm]{1}
  \quad\text{for every }x\in\X.
\]
\studentrecall{Each \emph{row} of the transition matrix is a probability distribution on $\Y$, so it sums to $1$.}
The channel is \keyword{discrete} because both alphabets are finite.
\end{definition}

\begin{definition}[memoryless]\label{def:memoryless}
The channel is \keyword{memoryless} if the output at each time depends only on
the input at \answerin[3.2cm]{that same time}, and not on earlier inputs or
outputs. Equivalently, for every $n$ and every $x^n\in\X^n$, $y^n\in\Y^n$,
\[
  p(y^n\mid x^n)=\prod_{i=1}^{n} \answerin[2.4cm]{p(y_i\mid x_i)}.
\]
\studentrecall{Only the current input matters, so the block law factorises into single-use laws $p(y_i\mid x_i)$.}
A channel that is both discrete and memoryless is called a
\keyword{discrete memoryless channel} (DMC).
\end{definition}

\begin{remark}
The product formula quietly assumes two things: that the channel does not
remember past inputs (memorylessness) and that the input $x_i$ is not allowed to
depend on the past outputs $y^{i-1}$ (no feedback). Both are needed for the
factorisation as written, and both hold in the setting of this lecture.
\end{remark}

\block{Capacity}

\begin{definition}[information channel capacity]\label{def:information-capacity}
The (``information'') \keyword{channel capacity} of a discrete memoryless
channel is
\[
  C=\max_{p(x)}\MI(X;Y),
\]
where the maximum is taken over all possible
\answerin[4.2cm]{input distributions $p(x)$}.
\studentrecall{Maximise over the input law $p(x)$; the transition matrix $p(y\mid x)$ is fixed.}
The joint law is always $p(x,y)=p(x)\,p(y\mid x)$, so choosing $p(x)$ chooses
both the output distribution $p(y)=\sum_xp(x)p(y\mid x)$ and the value of
$\MI(X;Y)$.
\end{definition}

\begin{remark}
Capacity also has an \keyword{operational} definition: the largest rate, in bits
per channel use, at which information can be sent with error probability tending
to $0$. Shannon's second theorem (the channel coding theorem) says that the
information capacity and the operational capacity coincide --- which is why the
word ``information'' is usually dropped. This lecture computes the
$\max_{p(x)}\MI(X;Y)$ side; the coding theorem is the next lecture's business.
\end{remark}

\block{The one move used in every example}

Because $\MI(X;Y)=\Ent(Y)-\Ent(Y\mid X)$ and
\[
  \Ent(Y\mid X)=\sum_{x}p(x)\,\Ent(Y\mid X=x),
\]
the second term is an average of \emph{fixed} row entropies of the transition
matrix. So maximising $\MI(X;Y)$ means making $\Ent(Y)$ large while keeping the
average row entropy small. When all rows have the \emph{same} entropy, the
second term does not depend on $p(x)$ at all and only $\Ent(Y)$ has to be
maximised --- this is the shortcut behind Sections \ref{sec:bsc} and
\ref{sec:symmetric}.

\begin{yourturn}
Write down the two-term decomposition of $\MI(X;Y)$ that this lecture uses, and
say which term the channel fixes and which term the input distribution moves.
\studentrecall{$\MI=\Ent(Y)-\Ent(Y\mid X)$; the rows of $p(y\mid x)$ fix the conditional entropies, while $p(x)$ moves their weights and $\Ent(Y)$.}
\ifshowanswers\else\workspace[2]\fi
\answerprose{$\MI(X;Y)=\Ent(Y)-\Ent(Y\mid X)$ with
$\Ent(Y\mid X)=\sum_xp(x)\Ent(Y\mid X=x)$. The channel fixes each row entropy
$\Ent(Y\mid X=x)$; the input distribution chooses their weights, and through
$p(y)=\sum_xp(x)p(y\mid x)$ it also controls $\Ent(Y)$.}
\end{yourturn}

\recall{What two assumptions are hidden in $p(y^n\mid x^n)=\prod_ip(y_i\mid x_i)$?}
